\section{Stromy}\label{sec:stromy}

Stromy tvoří speciální oblast teorie grafů. V~informatice jde o~jeden z~nejčastěji používaných typů grafů,~a~proto jim věnujeme samostatnou sekci,~v~níž si připomeneme některé základní pojmy.

\begin{definition}[Strom]\label{def:strom}
    Neorientovaný graf $G=(V,E)$ je strom,~pokud
    \begin{enumerate}[label=(\roman*)]
        \item je souvislý\footnote{Mezi každou dvojicí vrcholů vede cesta.}
        \item a~neobsahuje kružnici.
    \end{enumerate}
\end{definition}

Je dobré si uvědomit,~že mezi každou dvojicí vrcholů ve stromě vede \textbf{právě jedna} cesta. Zároveň z~definice stromu plyne,~že odebráním libovolné hrany se nutně poruší souvislost: přestane existovat cesta mezi koncovými vrcholy odebrané hrany.

\begin{figure}[h]
    \centering
    \begin{subfigure}{6cm}
        \centering
        \begin{tikzpicture}[scale=0.75]
    \foreach \name/\x/\y in {a/0/2,b/-1.7/0.7,c/1.7/0.7,d/-2.4/-0.8,e/-0.9/-0.8,f/1.2/-0.8,g/2.4/-0.8}
        \node[vertex,fill=black,minimum size=2mm] (\name) at (\x,\y) {};
    \foreach \a/\b in {a/b,a/c,b/d,b/e,c/f,c/g}\draw[edge](\a)--(\b);
\end{tikzpicture}

        \caption{Příklad korektního stromu.}
        \label{subfig:strom_priklad1}
    \end{subfigure}
    \qquad
    \begin{subfigure}{6cm}
        \centering
        \input{01-grafalgo/images/ch01_strom_priklad2.tex}
        \caption{Graf,~který není stromem,~protože obsahuje kružnici.}
        \label{subfig:strom_priklad2}
    \end{subfigure}
    \qquad
    \begin{subfigure}{6cm}
        \centering
        \begin{tikzpicture}[scale=0.72]
    \foreach \name/\x/\y in {a/-2/2,b/-3/0.5,c/-1/0.5,d/-3.6/-1,e/-2.4/-1,f/-1/-1,g/2/1.3,h/1/-0.4,i/3/-0.4}
        \node[vertex,fill=black,minimum size=2mm] (\name) at (\x,\y) {};
    \foreach \a/\b in {a/b,a/c,b/d,b/e,c/f,g/h,g/i}\draw[edge](\a)--(\b);
\end{tikzpicture}

        \caption{Graf,~který není stromem,~protože není souvislý.}
        \label{subfig:strom_priklad3}
    \end{subfigure}
\end{figure}

Poslední příklad \ref{subfig:strom_priklad3} sice není z~definice strom,~avšak čtenář by mohl namítnout,~že jednotlivé ``části'' daného grafu tvoří stromy. Tomuto typu grafů se říká (celkem pochopitelně) \emph{lesy},~neboť jejich \emph{komponenty souvislosti}\footnote{Intuitivně se jedná o~část grafu,~kde mezi každými dvěma vrcholy vede cesta. Formální zavedení termínu se provádí přes \emph{binární relaci}. Pro zájemce: \url{https://en.wikipedia.org/wiki/Component_(graph_theory)}.} tvoří stromy (viz obrázek \ref{fig:les}).

\begin{definition}[Les]\label{def:les}
    Graf $G=(V,E)$ je les,~pokud každá jeho komponenta souvislosti je strom.
\end{definition}

\begin{figure}[h]
    \centering
    \begin{tikzpicture}[scale=0.8]
    \foreach \name/\x/\y in {a/-3/1.8,b/-4/0.4,c/-2/0.4,d/-4.5/-1,e/-3.5/-1,f/-2/-1}
        \node[vertex,fill=black,minimum size=2mm] (\name) at (\x,\y) {};
    \foreach \a/\b in {a/b,a/c,b/d,b/e,c/f}
        \draw[edge] (\a)--(\b);

    \foreach \name/\x/\y in {
        u1/2.2/1.3,
        u2/3.325/0.65,
        u3/3.325/-0.65,
        u4/2.2/-1.3,
        u5/1.075/-0.65,
        u6/1.075/0.65
    }
        \node[vertex,fill=black,minimum size=2mm] (\name) at (\x,\y) {};
    \foreach \a/\b in {1/2,2/3,3/4,4/5,5/6}
        \draw[edge] (u\a)--(u\b);

    \node[vertex,fill=black,minimum size=2mm] at (5.3,-0.4) {};
\end{tikzpicture}

    \caption{Příklad lesa.}
    \label{fig:les}
\end{figure}

Ekvivalentně můžeme les definovat jako \emph{disjunktní sjednocení stromů}. Jako poslední si uveďme speciální typ stromu,~se kterým budeme dále nejčastěji pracovat,~a~to tzv. \emph{binární strom}. Blíže se s~ním setkáme v~sekci \smartref[][,~kdy budeme rozebírat tzv. \emph{binární haldu}]{\showrefs}{sec:halda}{o~haldě}.

\begin{definition}[Binární strom]\label{def:binarni_strom}
    Strom $T=(V,E)$ nazveme binární,~pokud
    \begin{enumerate}[label=(\roman*)]
        \item je zakořeněný
        \item a~každý vrchol má nejvýše dva syny.
    \end{enumerate}
\end{definition}

Jeden z~vrcholů stromu $T$ je označen jako \emph{kořen}. Typicky jej zakreslujeme nahoru nad všechny ostatní vrcholy a~syny každého vrcholu umisťujeme vždy níže oba na stejnou úroveň (viz obrázek \ref{fig:binarni_strom_priklad}).

\begin{figure}[h]
    \centering
    \begin{tikzpicture}[x=1cm,y=1cm]

% vrcholy
\node[vertex, label=above:{kořen}] (r) at (0,0) {};
\node[vertex, label=above left:{levý syn}]  (l) at (-1.8,-1.4) {};
\node[vertex, label=above right:{pravý syn}] (p) at (1.8,-1.4) {};

% plné hrany
\draw[edge] (r) -- (l);
\draw[edge] (r) -- (p);

% naznačení pokračování levého podstromu
\draw[edge, dashed] (l) -- ++(-0.9,-1.0);
\draw[edge, dashed] (l) -- ++(0.9,-1.0);

% naznačení pokračování pravého podstromu
\draw[edge, dashed] (p) -- ++(-0.9,-1.0);
\draw[edge, dashed] (p) -- ++(0.9,-1.0);

\end{tikzpicture}
    \caption{Příklad binárního stromu.}
    \label{fig:binarni_strom_priklad}    
\end{figure}
